\subsection{A Nivel de Miembro}

\subsubsection{Regla: orden de las etiquetas \texttt{param}, \texttt{returns} y \texttt{exception}}

\textbf{Regla:} En un método documentado por su complejidad, las etiquetas van en este orden: \texttt{<param>} por cada parámetro, luego \texttt{<returns>}, luego \texttt{<exception>} si aplica (Microsoft, 2026c).

\textbf{Código correcto:}
\begin{lstlisting}[style=csharpstyle]
/// <summary>
/// Calcula el puntaje final aplicando el multiplicador de la casilla actual.
/// </summary>
/// <param name="baseScore">Puntaje base acumulado antes del multiplicador.</param>
/// <param name="tile">Casilla en la que se detuvo el jugador.</param>
/// <returns>El puntaje final despues de aplicar el multiplicador de la casilla.</returns>
/// <exception cref="ArgumentOutOfRangeException">Si <paramref name="baseScore"/> es negativo.</exception>
public int CalculateFinalScore(int baseScore, BoardTile tile)
{
    if (baseScore < 0)
    {
        throw new ArgumentOutOfRangeException(nameof(baseScore));
    }

    return baseScore * tile.ScoreMultiplier;
}
\end{lstlisting}

\textbf{Código incorrecto:}
\begin{lstlisting}[style=csharpstyle]
/// <returns>El puntaje final despues de aplicar el multiplicador de la casilla.</returns>
/// <exception cref="ArgumentOutOfRangeException">Si <paramref name="baseScore"/> es negativo.</exception>
/// <param name="baseScore">Puntaje base acumulado antes del multiplicador.</param>
/// <param name="tile">Casilla en la que se detuvo el jugador.</param>
public int CalculateFinalScore(int baseScore, BoardTile tile)
{
    if (baseScore < 0)
    {
        throw new ArgumentOutOfRangeException(nameof(baseScore));
    }

    return baseScore * tile.ScoreMultiplier;
}
\end{lstlisting}
